\section{Harness 作为未来的服务模板}

\subsection{从 Service Template 到 Harness Template}

Birgitta 提出了一个引人深思的观察：大多数组织只有\textbf{两到三个主要技术栈}------并非每个应用都是独特的``雪花''（snowflake）。这让她设想了一个未来场景：团队可以从一组针对\textbf{常见应用拓扑}（common application topologies）的 harness 中进行选择来启动项目。

这个设想直接对应了当今软件工程中的 \textbf{service template}（服务模板）概念。Service template 帮助团队在``黄金路径''（golden path）上实例化新服务，提供标准化的项目结构、CI/CD 配置、监控集成等。

\begin{knowledgebox}{Service Template 的演进}
\begin{table}[H]
\centering
\begin{tabular}{lll}
\toprule
\textbf{维度} & \textbf{传统 Service Template} & \textbf{未来 Harness Template} \\
\midrule
包含内容 & 项目结构、配置、CI/CD & 自定义 linter、结构性测试、 \\
 & & 知识文档、上下文提供者 \\
目标用户 & 人类开发者 & AI agent + 人类 \\
初始化 & 一次性脚手架 & 持续演进的约束系统 \\
维护方式 & 手动更新、回馈上游 & agent 驱动的迭代改进 \\
\bottomrule
\end{tabular}
\end{table}
\end{knowledgebox}

\subsection{分叉与同步挑战}

Birgitta 敏锐地指出了一个潜在问题：在当前的 service template 实践中，团队基于模板起步后会根据自身需求进行定制，然后尝试将改进回馈给上游模板------但其他团队往往难以整合这些更新。这种\textbf{分叉与同步}的挑战在 harness 模式下是否会重现？

\begin{warningbox}{Harness 的分叉风险}
当每个团队都基于通用 harness 模板进行定制时，可能出现以下问题：
\begin{itemize}[leftmargin=1.5em]
    \item 不同团队的 harness 逐渐分化，失去统一标准的意义
    \item 将特定项目的改进合并回通用模板变得困难
    \item 通用模板的更新难以推送到已高度定制的项目 harness 中
    \item 团队可能陷入``harness 债务''------维护 harness 本身的负担越来越重
\end{itemize}
这与 service template 面临的问题本质相同，但由于 harness 包含更多动态组件（如 LLM 驱动的检查逻辑），分叉后的同步可能更加困难。
\end{warningbox}

\subsection{本章小结}

Harness 有可能成为 service template 的进化形态，为 AI 驱动的开发提供标准化的控制框架。但这也意味着它将继承 service template 的分叉与同步挑战，甚至可能因为 harness 的动态特性而变得更加复杂。组织需要在标准化与定制化之间找到平衡。

\newpage

